%% kapa: exam=1 round=36 year=2025 period=1 form=A subject=economics subject_ko=경제학원론 seq=15 number=55
%% source: Q-Net 감정평가사 기출문제 (cst00309) artlSeq=5231525 file=raw/1008/5231525/2025년 제36회 감정평가사 1차시험 1교시 문제지 원본.pdf
%% conversion: kordoc 4.24.51 -> tools/gen1.py (auto); review=proofread-sonnet; flags=none
\begin{question}{55}
갑과 을의 동시선택게임에서 전략이 각각 L, R과 T, D일 때, 다음 설명 중 옳은 것은? (단, 괄호 안 왼쪽은 갑, 오른쪽은 을의 보수)
\begin{center}
{\small
\begin{tabular}{|c|c|c|c|}
\hline
\multicolumn{2}{|c|}{} & \multicolumn{2}{c|}{을} \\ \cline{3-4}
\multicolumn{2}{|c|}{} & T & D \\ \hline
\multirow{2}{*}{갑} & L & (5, 5) & (0, 20) \\ \cline{2-4}
 & R & (20, 30) & (25, 0) \\ \hline
\end{tabular}}
\end{center}
\begin{choices}
  \item 갑의 전략 L에 대한 을의 최선반응(best response)은 T이다.
  \item 을의 전략 T에 대한 갑의 최선반응은 L이다.
  \item 갑은 우월전략이 없다.
  \item T는 을의 우월전략이다.
  \item 전략조합 (R, T)는 내쉬균형이다.
\end{choices}
\end{question}
\answer{5}
